\documentclass[11pt]{article}
\usepackage{amsmath,amsthm,amssymb}
\usepackage[hidelinks]{hyperref}

\newtheorem{theorem}{Theorem}
\newtheorem{lemma}{Lemma}
\newcommand{\chip}{\chi_{\mathrm p}}
\newcommand{\ind}{\mathbf{1}}
\newcommand{\file}[1]{\path{#1}}
\setlength{\emergencystretch}{2em}

\title{A Computer-Assisted Counterexample to the 22-Color Conjecture\protect\\
for Bridge-Connected $K_6$ Chains}
\author{Aurantium}
\date{}

\begin{document}
\maketitle

\begin{abstract}
Let $H_t$ consist of $t$ vertex-disjoint copies of $K_6$, with a single
bridge between each consecutive pair and distinct left and right bridge
vertices within every clique. We give a computer-assisted proof that
$\chip(H_t)\leq 23$ for every $t\geq 1$ and $\chip(H_t)=23$ for every
$t\geq 75$. In particular, the 450-vertex graph $H_{75}$ disproves
Conjecture~2 of Furma\'nczyk, G\"oz\"upek and \"Ozkan
(arXiv:2511.12761v1), and 23 is the sharp uniform upper bound for this
family. For the lower bound, every 22-coloring yields a necessary
scheduling relaxation. An integer decreasing-rank certificate with
16,646,562 records rules out a 75-row schedule; two checker
implementations have each checked all 1,755,492,299 labeled choices.
For the upper bound, an explicit 480-clique periodic coloring is
verified using graph distances. The proof uses exact integer checks
and does not require trusting the certificate generator. We do not
establish that 75 is the least counterexample length.
\end{abstract}

\section{The graph family and main result}

We work with finite, simple, undirected, unweighted graphs in the usual
ZFC foundations. A \emph{packing $k$-coloring} of a graph $G$ is a map
$c:V(G)\to\{1,\ldots,k\}$ such that distinct vertices of color $i$
have graph distance strictly greater than $i$. The packing chromatic
number $\chip(G)$ is the least such $k$. The numerical color labels
matter: an arbitrary permutation of them need not preserve a packing
coloring.

For each integer $t\geq 1$, let $H_t$ have pairwise disjoint vertex sets
\[
V_j=\{L_j,A_j,B_j,C_j,D_j,R_j\},\qquad 1\leq j\leq t.
\]
Each $V_j$ induces a $K_6$. The only edges between different cliques are
$R_jL_{j+1}$ for $1\leq j<t$. In particular, $L_j\ne R_j$, and
neighboring cliques do not share vertices. The proposed uniform
22-color bound is
\begin{equation*}
\begin{gathered}
\forall t\geq 1\quad \exists c:V(H_t)\longrightarrow\{1,\ldots,22\},\\
u\ne v,\quad c(u)=c(v)=i
\quad\Longrightarrow\quad d_{H_t}(u,v)>i.
\end{gathered}
\tag{T}\label{eq:target}
\end{equation*}
This is an upper-bound assertion; it does not assert that every finite
chain has packing chromatic number exactly 22.

The assertion comes from Conjecture~2 in Section~3.2.1 of
Furma\'nczyk, G\"oz\"upek and \"Ozkan~\cite[PDF p.~12]{FGO}.
Their Definition~2 partitions a path into vertex-disjoint blocks,
each identified with a path inside a copy of the prescribed graph
\cite[PDF p.~3]{FGO}. For $K_6$ and $2\leq\ell\leq 6$, the first and
last vertices of each $\ell$-vertex block are distinct. Name them
$L_j$ and $R_j$. All internal block edges are already clique edges,
and the only edges between consecutive blocks are the bridges
$R_jL_{j+1}$. Conversely, in every clique of $H_t$, choose any simple
$\ell$-vertex path from $L_j$ to $R_j$. This realizes the original
construction, so
\[
H_t\cong P_{\ell t}\diamond_{\ell}K_6
\qquad(2\leq\ell\leq 6).
\]
This is also the relevant specialization of Proposition~20
\cite[PDF p.~10]{FGO}. Thus the graph family here is the one in the
conjecture, with bridge-connected, vertex-disjoint cliques.

\begin{theorem}\label{thm:main}
For the graph family above,
\[
\chip(H_t)\leq 23\quad(t\geq 1),\qquad
\chip(H_t)=23\quad(t\geq 75).
\]
In particular, $\chip(H_{75})=23$, and 23 is the sharp uniform upper
bound for the family.
\end{theorem}

The graph $H_{75}$ has 450 vertices and 1199 edges. Its complete edge
list and a 23-coloring are supplied in
\file{certificates/H75_graph_and_23color_witness.json}.
The impossibility of a 22-coloring follows from the necessary
relaxation, rank induction, and complete finite certificate below.
It is not an inference from an unsuccessful periodic search or an
exclusion restricted to a paired-color template. We do not prove
that 75 is the smallest counterexample length or determine all
packing chromatic numbers for $t<75$.

\section{Boundary cases}

Since $H_1=K_6$, we have $\chip(H_1)=6$. The diameter of $H_2$ is 3.
Each color at least 3 can therefore occur at most once, while each
of colors 1 and 2 can occur at most twice. Nine colors can cover at
most $2+2+7=11$ vertices, fewer than the 12 vertices of $H_2$.
In the order $(L,A,B,C,D,R)$, the two rows
\[
(3,1,2,4,5,6),\qquad (7,1,2,8,9,10)
\]
give a packing 10-coloring. Consequently $\chip(H_2)=10$.

\section{Distances and a necessary scheduling relaxation}

\begin{lemma}[Exact distances]\label{lem:distance}
If $x\in V_j$, $y\in V_{j+r}$, and $r\geq 1$, then
\begin{equation*}
d(x,y)=2r-1+\ind[x\ne R_j]+\ind[y\ne L_{j+r}].
\tag{1}\label{eq:distance}
\end{equation*}
\end{lemma}

\begin{proof}
Every path between the two vertices must traverse $r$ bridges.
Within each of the $r-1$ strictly intermediate cliques, it must also
travel from the left bridge vertex to the distinct right bridge
vertex, using at least one internal edge. The initial and final
cliques contribute the two endpoint costs in
\eqref{eq:distance}. Since every clique is complete, these lower
bounds can all be attained simultaneously. Leaving this segment of
the bridge chain and returning cannot shorten a path.
\end{proof}

In particular, $d(x,y)\leq 2r+1$. Set
\[
q_i=\lceil i/2\rceil.
\]
If color $i$ occurs in two cliques whose indices differ by $r<q_i$,
then $2r+1\leq i$, contradicting the packing condition. Occurrences
of color $i$ must therefore have clique indices differing by at
least $q_i$. This is a necessary condition, not a sufficient one.

Every $K_6$ uses six distinct colors, at most two of which are 1 and
2. Suppose that a 22-coloring exists. In each clique, retain any
four colors greater than 2 and discard all other occurrences.
The result is a necessary schedule with $t$ rows: each row selects
four distinct tasks, corresponding to colors 3 through 22, and
successive selections of the task for color $i$ are separated by
at least $\lceil i/2\rceil$ rows. The waiting parameters of these
twenty tasks are
\begin{equation*}
2,2,3,3,\ldots,11,11.
\tag{2}\label{eq:parameters}
\end{equation*}
Deleting occurrences cannot reduce the gaps between retained
occurrences. This model also discards every endpoint restriction,
so it is weaker than the original coloring problem. Hence
\begin{equation*}
\begin{gathered}
H_t\text{ has a packing 22-coloring}\\
\Longrightarrow
\text{the tasks in \eqref{eq:parameters} admit a legal $t$-row schedule}\\
\text{with exactly four tasks selected in each row.}
\end{gathered}
\tag{3}\label{eq:relaxation}
\end{equation*}
This argument does not assume that colors 1 and 2 occur in every
clique. It also does not require colors $2q-1$ and $2q$ to occur
together. Neither restriction is used in the proof.

\section{Scheduling states and exact transitions}

Immediately before a row, a task of waiting parameter $q$ has a
residual waiting time $r\in\{0,\ldots,q-1\}$. It is available
exactly when $r=0$. A step selects four of the currently available
task copies. A selected task has new residual $q-1$; an unselected
task has new residual $\max(r-1,0)$. This rule expresses exactly
the requirement that selections of a task be separated by at
least $q$ rows.

The two tasks with the same parameter $q$ are interchangeable
within this necessary relaxation. We may therefore sort their
residuals as $0\leq a_q\leq b_q<q$. This sorting takes a quotient
by the symmetry of task labels; it does not force the tasks to be
selected together. Every labeled schedule projects to a sequence
of these sorted states. The full state domain is
\[
\mathcal S=\prod_{q=2}^{11}\{(a,b):0\leq a\leq b<q\},
\qquad
|\mathcal S|=\prod_{q=2}^{11}\frac{q(q+1)}2
=9\,336\,040\,560\,000.
\]
The initial state $s_0$ has all twenty residuals equal to zero,
because no selections precede the first row. The complete checkers
enumerate every labeled four-element subset of available task
copies at each listed state. Distinct labeled choices that become
the same sorted successor are still covered individually.

\section{Future-capacity inequalities}

\begin{lemma}[Capacity of a single task]\label{lem:capacity}
A task with parameter $q$ and current residual $r$ can be selected
at most $M_q(h,r)$ times in the next $h$ rows, where
\begin{equation*}
M_q(h,r)=
\begin{cases}
0,&h\leq r,\\
1+\lfloor(h-1-r)/q\rfloor,&h>r
\end{cases}
\tag{4}\label{eq:capacity}
\end{equation*}
\end{lemma}

\begin{proof}
Number the future rows $0,\ldots,h-1$. The earliest allowed row is
$r$, and subsequent selections must be at least $q$ rows apart.
Thus the earliest possible selection times are
$r,r+q,r+2q,\ldots$. Counting the terms at most $h-1$ gives
\eqref{eq:capacity}.
\end{proof}

For a state $s=((a_q,b_q))_{q=2}^{11}$, define
\[
C_h(s)=\sum_{q=2}^{11}\bigl(M_q(h,a_q)+M_q(h,b_q)\bigr).
\]
Four tasks are needed in each row. It follows that
\begin{equation*}
C_h(s)<4h\quad\Longrightarrow\quad
s\text{ has no legal extension of $h$ rows}.
\tag{5}\label{eq:obstruction}
\end{equation*}
This is an integer necessary inequality. It uses no numerical
optimization, and it need not characterize every feasible state.

\section{A finite decreasing-rank certificate}

\begin{lemma}[Decreasing ranks]\label{lem:rank}
Let $D\subseteq\mathcal S$ be finite with $s_0\in D$, and let
$R:D\to\mathbb Z_{>0}$. Suppose that, for every $s\in D$ and
every legal one-step successor $s'$, at least one of the following
holds:
\begin{enumerate}
\item[(A)] There is an integer $h$ with $1\leq h<R(s)$ such that
$C_h(s')<4h$.
\item[(B)] $s'\in D$ and $R(s')<R(s)$.
\end{enumerate}
Then no state $s\in D$ admits a legal schedule of length $R(s)$.
\end{lemma}

\begin{proof}
Use strong induction on the positive integer $R(s)$. Suppose a
legal schedule of length $R(s)$ starts at $s$. After its first
step, $R(s)-1$ rows remain. In case (A), the first $h$ rows of this
remaining schedule contradict \eqref{eq:obstruction}. In case (B),
the first $R(s')$ remaining rows contradict the induction
hypothesis. If there is no legal first step, the conclusion is
immediate.
\end{proof}

The lemma does not require $D$ to contain every possible state.
It also does not require trusting a search procedure: the required
check is that all legal transitions from every listed state are
covered by (A) or (B).

\section{The complete lower-bound certificate}

\subsection{Records and exact encoding}

The principal certificate is distributed as a separate compressed
asset, \file{H75_rank_certificate.bin.gz}. Its decompressed path
in the accompanying package is
\file{certificates/full_fast50m.bin}. It consists of 16,646,562
records, each an eight-byte little-endian unsigned integer, with
no header. Its total length is 133,172,496 bytes. Every recorded
rank is positive, and the initial state has rank 75. The SHA-256
of the decompressed certificate is
\begin{center}
\scriptsize\texttt{082c86b92fa216f25c8f4dcb2abf97bfa88c3f98ca93b2de150041e8eb5e4fb0}.
\end{center}
The compressed asset has 49,164,961 bytes and SHA-256
\begin{center}
\scriptsize\texttt{eb42565bf1c6e3e0e732583c13564886a05416c889c8b00dcedffd9ce6ad94b6}.
\end{center}
Both identities are also specified in
\file{CERTIFICATE_ASSET.json}.

For each $q$, put
\[
B_q=\frac{q(q+1)}2,\qquad
u_q=\frac{b_q(b_q+1)}2+a_q.
\]
The mixed-radix state identifier is
\[
\operatorname{id}(s)=
\sum_{q=2}^{11}u_q\prod_{v=2}^{q-1}B_v,
\]
where an empty product is 1. A 64-bit record $w$ encodes
\[
\mathrm{key}=w\bmod 2^{44}=\operatorname{id}(s)+1,
\qquad
R(s)=\left\lfloor w/2^{44}\right\rfloor.
\]
The records are ordered by strictly increasing key, and the
initial state has key 1. Every valid state identifier is less
than $2^{44}-1$. To decode a state, recover successive digits
$u_q$ by Euclidean division by the radices $B_q$. The unique
triangular interval
\[
\frac{b(b+1)}2\leq u_q<\frac{(b+1)(b+2)}2
\]
determines $b_q=b$, and then $a_q=u_q-b(b+1)/2$.
The full format and checker contract are given in
\file{docs/CERTIFICATE_FORMAT.md}.

\subsection{The checking algorithm}

The following pseudocode specifies a direct complete check. Here
$Q=(2,2,3,3,\ldots,11,11)$, and the residuals in each equal-$Q$
pair are kept in sorted order.
\begin{quote}\small
\begin{verbatim}
Read every little-endian record and validate its encoding,
positive rank, state domain, and strictly increasing key.
Store the resulting map R, and require R[all_zero] == 75.
For every state s in R:
    available = task copies with residual zero in s
    For every four-element subset A of available:
        t[i] = Q[i]-1 if i is in A
               else max(s[i]-1, 0)
        Sort each equal-Q pair in t.
        If 0 < R.get(t, 0) < R[s]:
            Accept this transition.
        Otherwise:
            Require some h in 1..min(72, R[s]-1) with
                sum(M(Q[i], h, t[i])) < 4*h.
Accept only if every transition passes.
\end{verbatim}
\end{quote}
The finite certificate metadata, including byte count, record
count and root rank, must match the stated certificate. Restricting
the capacity search to $h\leq72$ is sound: a single obstruction
with $1\leq h<R(s)$ suffices. There is no assumption that this
restricted search recognizes all impossible states. A transition
without such an obstruction must have a listed successor of
strictly smaller positive rank.

For efficient capacity checking, first decrease every residual to
$\max(r-1,0)$, obtaining a base state $b$. If an available task
with parameter $q$ is selected, its residual changes from 0 in
$b$ to $q-1$ in the successor. The resulting decrease in its
$h$-row capacity is exactly $\ind[q\nmid h]$. Thus, if the four
selected copies form $A$, the successor satisfies
\begin{equation*}
C_h(s')=C_h(b)-4+|\{a\in A:q_a\mid h\}|.
\tag{6}\label{eq:capacity-update}
\end{equation*}
This identity follows directly by counting possible selection
times $0,q,2q,\ldots$ and $q-1,2q-1,\ldots$. The implementations
can test capacity before doing the rank lookup. The original
checker also combines mathematically identical capacity tests;
neither optimization changes the obligation to cover every
labeled choice.

The package contains the original complete checker,
\file{code/verify_full_rank.cpp}; a portable copy,
\file{code/verify_full_rank_portable.cpp}; and a separately
implemented checker, \file{code/independent_rank_check.cpp}.
The portable copy replaces two compiler intrinsics with standard
integer loops, preserving the mathematical checks. The separate
implementation uses its own little-endian parsing, hash table,
and local successor encoding. Both implementations nevertheless
share the state model and capacity identity proved above; they
provide an implementation cross-check, not two proofs with no
shared mathematical assumptions.

Neither checker calls the generator, Z3, HiGHS, or another
optimization solver. All mathematical checks use integers;
floating point is used only for elapsed-time reporting. The
single-task capacity tables are constructed by explicitly
enumerating the earliest possible selection times, independently
of the generator's capacity routine.

\subsection{Complete verification counts and the lower bound}

Two retained complete executions, one of the portable checker
and one of the separate implementation, both passed with the
following counts. No restriction was imposed on the first step.
\begin{center}
\begin{tabular}{lr}
\hline
Quantity & Count\\
\hline
Verified states & 16,646,562\\
Labeled four-task choices & 1,755,492,299\\
Choices covered by capacity & 1,694,728,516\\
Choices covered by decreasing rank & 60,763,783\\
Initial rank & 75\\
Maximum recorded rank & 75\\
Largest capacity horizon used & 70\\
\hline
\end{tabular}
\end{center}
The capacity and rank counts sum to the total number of choices.
For every transition not excluded by capacity, the successor was
encoded and its smaller positive rank checked. The complete
mathematical outputs, with timing fields omitted, and the hashes
of the checker source files are retained in
\file{verification/negative_full_results.json}.
These are complete checks, rather than sampling results or
solver infeasibility logs.

Lemma~\ref{lem:rank}, together with $R(s_0)=75$, rules out a
75-row schedule in the necessary relaxation.
By \eqref{eq:relaxation}, $H_{75}$ has no packing 22-coloring.
Therefore
\begin{equation*}
\chip(H_{75})\geq23.
\tag{7}\label{eq:lower}
\end{equation*}

\section{An explicit periodic 23-color upper bound}

The file
\file{certificates/k6_23color_period480_certificate.json}
contains the positive certificate: 480 rows, with the six actual vertex colors in each row
ordered as $(L,A,B,C,D,R)$. All entries are in
$\{1,\ldots,23\}$. Its SHA-256 is
\begin{center}
\scriptsize\texttt{5e9136a6030a0827179c08d0a6b87efa1ec3606aabb561c429ad99068a337f31}.
\end{center}
Repeat the 480 rows periodically along the infinite chain. Each
row has six distinct colors.

The checker \file{code/verify_periodic_certificate.py}
uses the Python standard library. It constructs the adjacency list of 492 consecutive copies of $K_6$. It performs
breadth-first search (BFS) from each of the 2880 vertices in the
first period. It checks all vertex pairs whose later clique is
between 1 and 12 cliques after the source clique, and compares
the BFS distances with \eqref{eq:distance}. The completed check
reports
\begin{center}
\begin{tabular}{lr}
\hline
Quantity & Value\\
\hline
Status & PASS\\
Period & 480\\
Number of colors & 23\\
BFS sources & 2880\\
BFS/formula comparisons & 207,360\\
Same-color pairs in the BFS window & 17,536\\
Successive cyclic same-color pairs & 2880\\
Far-pair distance lower bound & 25\\
\hline
\end{tabular}
\end{center}
The full output is retained in
\file{outputs/period480_independent_check.json}.

This is a repeated straight chain. No edge is added from the
last clique of a period back to its first clique. Pairs crossing
a period boundary are checked in the explicitly constructed
following cliques. For any pair in the infinite chain, translate
the earlier clique by a multiple of 480 into the first period.
If the clique-index difference is at most 12, the translated
pair is among the checked pairs. If it is at least 13, then
\eqref{eq:distance} gives
\[
d(x,y)\geq 2\cdot13-1=25>23,
\]
so every such pair automatically satisfies the packing
condition. Pairs in a single clique are covered by the six
distinct colors in every row. This proves that the finite
window covers every possible conflict in the repeated coloring.

As an additional check, for each color $i$ the minimum actual
graph distance between consecutive cyclic occurrences is
exactly $i+1$, including occurrences across a period boundary.
For a fixed earlier vertex, increasing the target clique index
by one increases the distance in \eqref{eq:distance} by at
least one. Thus nonconsecutive occurrences of the same color
cannot be closer. These per-color checks are additional to,
and do not replace, the complete-window BFS check.

The periodically repeated rows therefore give a packing
23-coloring of the infinite chain. Restricting to its first $t$
cliques gives a packing coloring of every finite $H_t$, and hence
\begin{equation*}
\chip(H_t)\leq23\qquad(t\geq1).
\tag{8}\label{eq:upper}
\end{equation*}

The first 75 rows and the complete edge list of $H_{75}$ are also
stored in \file{certificates/H75_graph_and_23color_witness.json}.
The separate checker \file{code/verify_H75.py} confirms its 450
vertices and 1199 edges and performs 450 BFS runs. It checks all
101,025 unordered vertex pairs, including 8485 same-color pairs;
all pass. This confirms the finite positive witness. The
22-color lower bound is established separately by the negative
rank certificate.

\begin{proof}[Proof of Theorem~\ref{thm:main}]
Combining \eqref{eq:lower} and \eqref{eq:upper} gives
$\chip(H_{75})=23$. For $t\geq75$, the first 75 cliques induce
an isometric copy of $H_{75}$: the bridge chain offers no shorter
path through later cliques. A packing 22-coloring of $H_t$
would therefore restrict to one of $H_{75}$, a contradiction.
The uniform upper bound \eqref{eq:upper} proves
$\chip(H_t)=23$ for every $t\geq75$. Because a member of the
family requires 23 colors, this uniform constant is sharp.
\end{proof}

\section{Reproduction and verification scope}

The proof depends on the mathematical reduction, the rank
induction and the complete negative-certificate checks, together
with the periodic positive certificate on actual vertex
positions. The generator and its search decisions are not proof
dependencies once the exact certificate is supplied. The
additional Python sample checker is supplementary and cannot
replace complete negative verification. File hashes identify
the data being checked; they do not replace the mathematical
checks.

Reproduction requires a 64-bit environment, Python 3.10 or later,
and a C++17 compiler. The checkers use standard libraries and
require no SAT or mixed-integer optimization solver. The full
certificate and lookup table require approximately 400 MB of
memory, before environment-dependent overhead. The accompanying
\file{REPRODUCIBILITY.md} gives integrity and decompression
commands, compiler-specific commands, and expected verification
counts. Download the compressed certificate asset matching the
source package, verify it using \file{CERTIFICATE_ASSET.json},
and decompress it to \file{certificates/full_fast50m.bin}.

With a GCC- or Clang-compatible compiler available, the principal
reproduction command from the package root is
\begin{quote}\small
\begin{verbatim}
python reproduce_primary.py
\end{verbatim}
\end{quote}
The wrapper checks the eight primary file hashes, compiles the
original complete checker, verifies the entire negative
certificate, checks the periodic and finite positive witnesses
by BFS, and runs the supplementary Python sample check. It
requires the stated complete counts and rejects a restricted
first step. It creates a fresh \file{recheck/} directory; a
repeat run can use a new \file{--out-dir}. Its compiler option
uses GCC/Clang command-line syntax, not MSVC syntax.

The second complete implementation can be run separately:
\begin{quote}\small
\begin{verbatim}
g++ -O3 -std=c++17 code/independent_rank_check.cpp \
    -o independent_rank_check
./independent_rank_check certificates/full_fast50m.bin
\end{verbatim}
\end{quote}
For MSVC, the reproduction instructions give commands for both
the separate implementation and the portable first checker.
Each normal complete run must report success, initial rank 75,
16,646,562 states, and 1,755,492,299 labeled choices, with the
capacity and rank counts stated above. The original checker's
\file{--first-high} option belongs to a different exploratory
certificate and must not be used for this theorem. The separate
checker's \file{formal_metadata: true} field denotes its fixed
certificate mode; it does not denote proof-assistant
formalization. Its self-test and fixture modes likewise do not
establish the theorem.

The positive checks can be run independently:
\begin{quote}\small
\begin{verbatim}
python code/verify_periodic_certificate.py \
    certificates/k6_23color_period480_certificate.json
python code/verify_H75.py \
    certificates/H75_graph_and_23color_witness.json
\end{verbatim}
\end{quote}
These establish upper bounds; the negative certificate is
separately required for the lower bound.

The summaries in \file{verification/negative_full_results.json}
and \file{outputs/} retain results of completed executions.
Preparing this manuscript does not itself constitute another
complete negative-certificate run. Reproduction should check the
actual certificate instead of treating a recorded PASS as a
proof by itself. As in an ordinary computer-assisted proof, the
checker implementations, their compilers and execution
environments remain within the trust boundary. Hashes do not
certify compiler semantics or the authenticity of unsigned
execution records. No proof-assistant kernel formalization is
provided, and no external expert-review endorsement is claimed.

The result does not determine the least counterexample length or
all values for shorter chains. The computations for $t\leq48$
reported in the original paper~\cite[Section~3.2.1, PDF p.~12]{FGO}
are consistent with a counterexample at $t=75$; we make no
assessment of unavailable original experimental data. No claim
of worldwide priority is made.

\section*{Acknowledgment of AI assistance}

This work was completed with AI assistance. Its mathematical
claims rest on the arguments, exact certificates and reproducible
checks presented here, rather than on generated text itself.

\begin{thebibliography}{1}
\bibitem{FGO}
Hanna Furma\'nczyk, Didem G\"oz\"upek, and Sibel \"Ozkan.
\newblock \emph{Packing coloring of graphs with long paths}.
\newblock arXiv:2511.12761v1, 16 November 2025.
\newblock \url{https://arxiv.org/abs/2511.12761v1}.
\newblock Definition~2: PDF p.~3; Proposition~20: PDF p.~10;
Section~3.2.1 and Conjecture~2: PDF p.~12.
\end{thebibliography}

\end{document}
